\documentclass[11pt,a4paper]{article}

\usepackage[utf8]{inputenc}
\usepackage[T1]{fontenc}
\usepackage{lmodern}
\usepackage[spanish,es-tabla,es-noshorthands]{babel}
\usepackage[a4paper,margin=2.3cm]{geometry}
\usepackage{xcolor}
\usepackage{listings}
\usepackage{tcolorbox}
\usepackage{enumitem}
\usepackage{titlesec}
\usepackage{parskip}
\usepackage{hyperref}
\usepackage{array}
\usepackage{colortbl}

\definecolor{azulOscuro}{HTML}{1E3C78}
\definecolor{azulMedio}{HTML}{2A5BB1}
\definecolor{grisCodigo}{HTML}{F5F6FA}
\definecolor{grisBorde}{HTML}{D0D4DD}
\definecolor{verdeOK}{HTML}{2D7D46}
\definecolor{rojoErr}{HTML}{B23B3B}
\definecolor{ambar}{HTML}{B07C00}

\hypersetup{
  colorlinks=true,
  linkcolor=azulMedio,
  urlcolor=azulMedio,
  pdftitle={Guia de pruebas - Fase 2}
}

\titleformat{\section}
  {\color{azulOscuro}\Large\bfseries}
  {\thesection.}{0.6em}{}
\titleformat{\subsection}
  {\color{azulMedio}\large\bfseries}
  {\thesubsection.}{0.5em}{}
\titleformat{\subsubsection}
  {\color{azulMedio}\normalsize\bfseries}
  {\thesubsubsection.}{0.5em}{}

\lstdefinestyle{src}{
  backgroundcolor=\color{grisCodigo},
  basicstyle=\ttfamily\small,
  frame=single,
  framerule=0.4pt,
  rulecolor=\color{grisBorde},
  numbers=none,
  showstringspaces=false,
  breaklines=true,
  xleftmargin=8pt,
  xrightmargin=8pt,
  aboveskip=6pt,
  belowskip=6pt,
}
\lstset{style=src}

\newtcolorbox{caja}[2][azulMedio]{
  colback=grisCodigo,
  colframe=#1,
  title=\textbf{#2},
  fonttitle=\color{white},
  boxrule=0.5pt,
  arc=2pt,
  left=6pt,right=6pt,top=4pt,bottom=4pt,
  before skip=8pt, after skip=8pt,
}

\newcommand{\test}[1]{%
  \subsection*{\color{azulOscuro}#1}
  \addcontentsline{toc}{subsection}{#1}
}

\title{\vspace{-1.2cm}\color{azulOscuro}\textbf{Compiladores 2026 --- Fase 2}\\[0.3em]
\large\color{black}Guía de pruebas en vivo}
\author{}
\date{}

\begin{document}
\maketitle
\vspace{-1.4cm}

\section{Setup}

\begin{caja}{Tkinter (recomendado para presentar)}
\begin{lstlisting}
python "files (1)/main.py"
\end{lstlisting}
\begin{itemize}[leftmargin=*,nosep]
  \item \textbf{F5} = Analizar.
  \item Botón \textbf{Ejemplos} (abajo a la izquierda) carga los 7
        programas precargados.
\end{itemize}
\end{caja}

\begin{caja}{Web (alternativa, si tenés XAMPP/WAMP corriendo)}
\begin{enumerate}[leftmargin=*,nosep]
  \item Copiar \texttt{files (1)/} a \texttt{htdocs/minicompiler/}.
  \item Abrir \url{http://localhost/minicompiler/index.php}.
  \item Botón \textbf{Analizar} arriba.
\end{enumerate}
\end{caja}

\section{Pruebas a presentar (orden recomendado)}

\subsection*{Test 1 --- Pipeline completo: frontend $\rightarrow$ TAC $\rightarrow$ backend optimizado}

\textbf{Código:}
\begin{lstlisting}
int a = 3;
int b = 4;
int c = a + b;
print(c);
\end{lstlisting}

\textbf{Qué mostrar:}
\begin{itemize}[leftmargin=*,nosep]
  \item Tab \textbf{Tokens}: 16 tokens detectados.
  \item Tab \textbf{Árbol}: AST con nodos
        \texttt{Program $\rightarrow$ Declaration $\rightarrow$ BinaryOp(+)}.
  \item Tab \textbf{Código Intermedio $\rightarrow$ Original}:
        \texttt{\$t1 = a + b; c = \$t1; print c}.
  \item Tab \textbf{Código Intermedio $\rightarrow$ Optimizado}:
        \texttt{c = 7; print 7} (reducción $\sim$30\,\%).
\end{itemize}

\textbf{Qué decir:} ``Acá se ve el flujo completo: el frontend produce el
TAC, el backend lo simplifica sin tocar el lenguaje fuente. Esa es la
frontera.''

\subsection*{Test 2 --- Constant Folding}

\textbf{Código:}
\begin{lstlisting}
int x = (3 + 4) * 2 - 5;
print(x);
\end{lstlisting}

\textbf{Esperado:}
\begin{itemize}[leftmargin=*,nosep]
  \item Original: 5 cuádruplos con \texttt{\$t1, \$t2, \$t3}.
  \item Optimizado: \texttt{x = 9; print 9}.
\end{itemize}

\textbf{Qué decir:} ``Como los dos operandos son constantes, el
optimizador evalúa la expresión en compilación. La aritmética desaparece
del código generado.''

\subsection*{Test 3 --- Identidades algebraicas}

\textbf{Código:}
\begin{lstlisting}
int x = 7;
int y = x + 0;
int z = x * 1;
int w = x * 0;
print(y); print(z); print(w);
\end{lstlisting}

\textbf{Esperado:}
\begin{itemize}[leftmargin=*,nosep]
  \item \texttt{y} y \texttt{z} se reducen a copias directas de \texttt{x}.
  \item \texttt{w} se reemplaza por \texttt{0} literal.
  \item Los temporales intermedios desaparecen por dead-code.
\end{itemize}

\textbf{Qué decir:} ``\texttt{x+0}, \texttt{x*1} y \texttt{x*0} se
reescriben sin necesidad de evaluar --- son identidades.''

\subsection*{Test 4 --- Branch Pruning}

\textbf{Código:}
\begin{lstlisting}
if (true) {
  print(1);
} else {
  print(2);
}
\end{lstlisting}

\textbf{Esperado:}
\begin{itemize}[leftmargin=*,nosep]
  \item Original: tiene \texttt{ifFalse true goto \$L1}, ambas ramas, etiquetas.
  \item Optimizado: solo \texttt{print 1}. La rama \texttt{else} y los
        saltos desaparecen.
\end{itemize}

\textbf{Qué decir:} ``Después de copy-propagation la condición queda como
literal \texttt{true}. La pasada de Branch Pruning detecta el salto
inútil y lo elimina; después Jump Threading limpia las etiquetas
huérfanas.''

\subsection*{Test 5 --- Control de flujo (while)}

\textbf{Código:}
\begin{lstlisting}
int n = 0;
int i = 0;
while (i < 3) {
  n = n + i;
  i = i + 1;
}
print(n);
\end{lstlisting}

\textbf{Qué mostrar:}
\begin{itemize}[leftmargin=*,nosep]
  \item En el TAC se ve la estructura del while:
        \texttt{\$L1: \dots ifFalse \dots goto \$L2 \dots goto \$L1 \dots \$L2:}.
  \item Las variables del lazo \textbf{no} se foldean (porque su valor
        cambia en cada iteración).
  \item El optimizador respeta los bordes de bloque (las etiquetas
        invalidan el entorno de propagación).
\end{itemize}

\textbf{Qué decir:} ``Acá se ve por qué la propagación se invalida en
saltos: si el optimizador asumiera que \texttt{i} sigue valiendo \texttt{0}
dentro del lazo, rompería el programa.''

\subsection*{Test 6 --- For completo}

\textbf{Código:}
\begin{lstlisting}
int s = 0;
for (int i = 0; i < 5; i = i + 1) {
  s = s + i;
}
print(s);
\end{lstlisting}

\textbf{Qué mostrar:} misma idea que el while; se ve cómo \texttt{for} se
traduce a TAC con etiqueta de inicio, condición, cuerpo, actualización y
salto.

\subsection*{Test 7 --- Recuperación ante errores}

\textbf{Código:}
\begin{lstlisting}
int x = 5@;
int x = 10;
string s = "sin cerrar
\end{lstlisting}

\textbf{Qué mostrar:}
\begin{itemize}[leftmargin=*,nosep]
  \item Tab \textbf{Errores} lista los 3 problemas:
  \begin{itemize}[leftmargin=*,nosep]
    \item \texttt{Carácter ilegal '@'} (léxico)
    \item \texttt{'x' ya fue declarado} (semántico)
    \item \texttt{Carácter ilegal '"'} (cadena sin cerrar)
  \end{itemize}
  \item A pesar de los errores, el TAC \textbf{se sigue generando} para
        la parte válida.
\end{itemize}

\textbf{Qué decir:} ``El compilador no aborta al primer error: reporta
todos y sigue generando lo que puede.''

\subsection*{Test 8 --- Unary minus (bug recién arreglado)}

\textbf{Código:}
\begin{lstlisting}
int x = -5;
int y = -x;
print(y);
\end{lstlisting}

\textbf{Esperado:}
\begin{itemize}[leftmargin=*,nosep]
  \item Original:
        \texttt{\$t1 = 0 - 5; x = \$t1; \$t2 = 0 - x; y = \$t2; print y}
  \item Optimizado: \texttt{x = -5; y = 5; print 5}
\end{itemize}

\textbf{Qué decir:} ``El unario menos se traduce a \texttt{0 - operando},
y el optimizador lo dobla automáticamente.''

\subsection*{Test 9 --- Coloreado del editor (solo Tkinter)}

\textbf{Código:}
\begin{lstlisting}
string saludo = "Hola Mundo";
int edad = 25;
boolean activo = true;
float pi = 3.1416;
\end{lstlisting}

\textbf{Qué mostrar:} después de F5, en el editor:
\begin{itemize}[leftmargin=*,nosep]
  \item \texttt{string, int, boolean, float} en celeste (palabras clave).
  \item \texttt{"Hola Mundo"} en verde --- \textbf{incluyendo las dos comillas}.
  \item \texttt{25}, \texttt{3.1416} en naranja.
  \item \texttt{true} en amarillo.
\end{itemize}

\textbf{Qué decir (si alguien pregunta):} ``Antes el highlight se cortaba
un par de caracteres antes del final del string; ahora usa el largo del
texto fuente, no el del valor procesado.''

\subsection*{Test 10 --- Vista visual del optimizador (la pieza más vendible)}

\textbf{Código (programa que produce mucha reducción):}
\begin{lstlisting}
int a = 3;
int b = 4;
int c = a + b;
int d = c * 2;
int e = d - 0;
int f = e * 1;
int g = f * 0;
print(g);
\end{lstlisting}

\textbf{Qué mostrar en la pestaña Código Intermedio:}
\begin{enumerate}[leftmargin=*,nosep]
  \item \textbf{Sub-tab Original} --- 14+ cuádruplos.
  \item \textbf{Sub-tab Optimizado} --- 1-2 cuádruplos (\texttt{print 0}).
  \item \textbf{Sub-tab Comparación} --- los dos lado a lado.
  \item \textbf{Sub-tab Info} --- explicación de las pasadas.
  \item \textbf{Métricas en el header}:
        \texttt{Cuádruplos: 14 $\rightarrow$ 2 $\cdot$ Reducción: 86\,\% $\cdot$ Temporales: 4 $\rightarrow$ 0}.
  \item \textbf{Tabla de traza} abajo: cada pasada que hizo cambios, con
        \texttt{iter / pasada / antes / después / delta}.
\end{enumerate}

\textbf{Qué decir:} ``Esta vista sirve también como apoyo didáctico: no
solo qué se redujo, sino \emph{por qué pasada} y \emph{en qué
iteración}.''

\section{Si algo falla en vivo}

\renewcommand{\arraystretch}{1.3}
\begin{tabular}{|p{4.5cm}|p{4cm}|p{6cm}|}
\hline
\rowcolor{grisCodigo}
\textbf{Síntoma} & \textbf{Causa probable} & \textbf{Salida rápida} \\
\hline
Web carga pero TAC sale vacío
  & PHP no encuentra Python
  & Caer al modo JS local (deshabilitar backend con devtools) o usar
    Tkinter \\
\hline
Tkinter no abre el menú de Ejemplos
  & Foco perdido en otra ventana
  & Click en el editor primero, después en \textbf{Ejemplos} \\
\hline
\texttt{python3: command not found} (web Linux)
  & Falta alias
  & \texttt{sudo ln -s /usr/bin/python /usr/bin/python3} \\
\hline
El árbol no se dibuja (web)
  & Tokens vacíos por error léxico
  & Mirar el tab \textbf{Errores}; corregir el código fuente \\
\hline
\end{tabular}

\section{Tips para la presentación}

\begin{itemize}[leftmargin=*]
  \item \textbf{Tené las 10 muestras pegadas en un archivo de texto
        aparte} --- copiar/pegar es más rápido y profesional que tipear
        en vivo.
  \item \textbf{Empezá por el Test 1} (pipeline completo) y \textbf{terminá
        con el Test 10} (vista visual con mucha reducción) --- el efecto
        \emph{wow} está en ver \texttt{14 $\rightarrow$ 2 cuádruplos}.
  \item Si presenta otra persona y se traba, los 7 ejemplos del menú
        \textbf{Ejemplos} ya están precargados como red de seguridad ---
        el \#4 (``Expresiones Complejas'') da \texttt{33 $\rightarrow$ 15}
        ($\sim$55\,\% reducción), buen fallback.
  \item Antes de la presentación: corré los 7 ejemplos una vez en orden
        para confirmar que la máquina anda fina.
\end{itemize}

\end{document}
